% GENERATED FILE. Edit canonical lessons, then run npm run generate:books.
% canonical-source-hash: fnv1a64:f5215d50678b489c
% canonical-lessons: FR-C52-poitrine, FR-C52-estomac, FR-C52-cheville, FR-C52-coude, FR-C52-poignet

\chapter{Chest, Stomach, Ankle, Elbow, Wrist}
\label{ch:fr-chest-stomach-ankle-elbow-wrist}
\hlchaptermodality{52}

\begin{chapteropening}
\textbf{By the end of this chapter:} \emph{I can name the chest, the stomach, an ankle, an elbow, a wrist.}

Complete the last lesson of chapter 52: i can name the chest, the stomach, an ankle, an elbow, a wrist.
\end{chapteropening}

\section[la poitrine]{la poitrine — the chest}

\label{lesson:FR-C52-poitrine}



\begin{quote}
Before the new one: say the French for the throat, then the French for a cheek.
\end{quote}

\subsection*{You'll want to know: la poitrine}
\textbf{la poitrine} — \textquotedblleft{}the chest\textquotedblright{}.

\begin{grammarlens}[title={What you've built}]
One more word for this chapter.
\end{grammarlens}

\subsection*{Your turn}
\begin{itemize}
\raggedright
  \item \emph{Say it:} \emph{la poitrine}
  \item \emph{Say it:} \emph{la poitrine}, once more
  \item \emph{Say it:} say \emph{la poitrine}
  \item \emph{Recall:} say the French for the throat, then the French for a cheek, then say \emph{la poitrine} again
\end{itemize}

\begin{tcolorbox}[breakable,colback=teal!4,colframe=teal!35!black,title={Before you move on}]
What does la poitrine mean? (\textquotedblleft{}The chest\textquotedblright{}.) Say it once more.
\end{tcolorbox}
\section[l'estomac]{l'estomac — the stomach}

\label{lesson:FR-C52-estomac}



\begin{quote}
Before the new one: say the French for a cheek, then the French for the chest.
\end{quote}

\subsection*{You'll want to know: l'estomac}
\textbf{l'estomac} — \textquotedblleft{}the stomach\textquotedblright{}.

\begin{grammarlens}[title={What you've built}]
One more word for this chapter.
\end{grammarlens}

\subsection*{Your turn}
\begin{itemize}
\raggedright
  \item \emph{Say it:} \emph{l'estomac}
  \item \emph{Say it:} \emph{l'estomac}, once more
  \item \emph{Say it:} say \emph{l'estomac}
  \item \emph{Recall:} say the French for a cheek, then the French for the chest, then say \emph{l'estomac} again
\end{itemize}

\begin{tcolorbox}[breakable,colback=teal!4,colframe=teal!35!black,title={Before you move on}]
What does l'estomac mean? (\textquotedblleft{}The stomach\textquotedblright{}.) Say it once more.
\end{tcolorbox}
\section[la cheville]{la cheville — an ankle}

\label{lesson:FR-C52-cheville}



\begin{quote}
Before the new one: say the French for the chest, then the French for the stomach.
\end{quote}

\subsection*{You'll want to know: la cheville}
\textbf{la cheville} — \textquotedblleft{}an ankle\textquotedblright{}.

\begin{grammarlens}[title={What you've built}]
One more word for this chapter.
\end{grammarlens}

\subsection*{Your turn}
\begin{itemize}
\raggedright
  \item \emph{Say it:} \emph{la cheville}
  \item \emph{Say it:} \emph{la cheville}, once more
  \item \emph{Say it:} say \emph{la cheville}
  \item \emph{Recall:} say the French for the chest, then the French for the stomach, then say \emph{la cheville} again
\end{itemize}

\begin{tcolorbox}[breakable,colback=teal!4,colframe=teal!35!black,title={Before you move on}]
What does la cheville mean? (\textquotedblleft{}An ankle\textquotedblright{}.) Say it once more.
\end{tcolorbox}
\section[le coude]{le coude — an elbow}

\label{lesson:FR-C52-coude}



\begin{quote}
Before the new one: say the French for the stomach, then the French for an ankle.
\end{quote}

\subsection*{You'll want to know: le coude}
\textbf{le coude} — \textquotedblleft{}an elbow\textquotedblright{}.

\begin{grammarlens}[title={What you've built}]
One more word for this chapter.
\end{grammarlens}

\subsection*{Your turn}
\begin{itemize}
\raggedright
  \item \emph{Say it:} \emph{le coude}
  \item \emph{Say it:} \emph{le coude}, once more
  \item \emph{Say it:} say \emph{le coude}
  \item \emph{Recall:} say the French for the stomach, then the French for an ankle, then say \emph{le coude} again
\end{itemize}

\begin{tcolorbox}[breakable,colback=teal!4,colframe=teal!35!black,title={Before you move on}]
What does le coude mean? (\textquotedblleft{}An elbow\textquotedblright{}.) Say it once more.
\end{tcolorbox}
\section[le poignet]{le poignet — a wrist}

\label{lesson:FR-C52-poignet}



\begin{quote}
Before the new one: say the French for an ankle, then the French for an elbow.
\end{quote}

\subsection*{You'll want to know: le poignet}
\textbf{le poignet} — \textquotedblleft{}a wrist\textquotedblright{}.

\begin{grammarlens}[title={What you've built}]
One more word for this chapter.
\end{grammarlens}

\subsection*{Your turn}
\begin{itemize}
\raggedright
  \item \emph{Say it:} \emph{le poignet}
  \item \emph{Say it:} \emph{le poignet}, once more
  \item \emph{Say it:} say \emph{le poignet}
  \item \emph{Recall:} say the French for an ankle, then the French for an elbow, then say \emph{le poignet} again
\end{itemize}

\begin{tcolorbox}[breakable,colback=teal!4,colframe=teal!35!black,title={Before you move on}]
What does le poignet mean? (\textquotedblleft{}A wrist\textquotedblright{}.) Say it once more.
\end{tcolorbox}
